\documentclass[../../../include/open-logic-section]{subfiles}

\begin{document}

\olfileid{sth}{z}{nat}
\olsection{选择我们的自然数}

%Let us take it for granted that the informal principle \stagesinf{}
%can be justified. It is also, though, worth commenting on the
%particular axiom we extracted from that informal principle.

在 \olref[infinity-again]{defnomega} 节，我们明确定义了“自然数”这一说法。该如何理解这一规定？它并非形而上学的断言，而仅仅是一个\emph{将}某些集合\emph{视为}自然数的决定。我们在 \olref[sfr][rel][ref]{sec}、\olref[sfr][arith][ref]{sec} 和 \olref[sfr][infinite][dedekindsproof]{sec} 等节中已触及这样做的理由。但我们还可以让这些理由更具针对性。

我们的无穷公理沿用了 \citet{VonNeumann1925}的思路。但这里还有另一条公理，我们本可以转而采用它：

\begin{defish}
\emph{Zermelo\citeyear{Zermelo1908Untersuchungen} 的无穷公理。} 存在一个集合 $A$，使得 $\emptyset \in A$ 且 $(\forall x \in A)\{x\} \in A$。
\end{defish}

假若我们采用的是 Zermelo 的公理，而非我们（受 von Neumann 启发）的无穷公理，我们同样会得到一个 Dedekind无限集，进而得到一个Dedekind代数。在 Zermelo 的方法中，该代数的特异元同样会是 $\emptyset$（我们用以代表 $0$ 的替代物），但单射将由映射 $x \mapsto \{x\}$ 给出，而非 $x \mapsto x \cup \{x\}$。这样做最直接的结果是，Zermelo 将 $2$ 视为 $\{\{\emptyset\}\}$，而我们（与 von Neumann 一道）将 $2$ 视为 $\{\emptyset, \{\emptyset\}\}$。

为何选择这条无穷公理而非另一条？主要的实际考量在于，von Neumann 的方法能很好地“向上扩展”以处理超限数。我们将从 \olref[ordinals][]{chap} 章开始探讨这一点。然而，仅从\emph{做算术}的简单视角来看，两种方法同样有效。所以，如果有人告诉你自然数\emph{就是}集合，那么一个明显的问题是：\emph{它们究竟是哪些集合？}

这个确切的问题因 \citet{Benacerraf1965}（贝纳塞拉夫）而闻名。但值得强调的是，它仅仅是我们早已多次遇到的现象中最著名的一例。基本要点如下：集合论为我们提供了一种\emph{模拟}一系列“直观”实体的方法——不仅包括实数、有理数、整数和自然数，也包括序对、函数和关系。然而，集合论从未为我们提供一种\emph{唯一的}模拟选择。\emph{总是}存在一些替代方案——直截了当地说——它们同样能很好地胜任。

\end{document}